\documentclass[11pt]{article}

\usepackage{amsmath, amssymb, amsthm}
\usepackage{geometry}
\usepackage{hyperref}
\usepackage{xcolor}
\usepackage{inputenc}
\geometry{margin=1in}

\newcommand{\R}{\mathbb{R}}

\title{Softening Duffing oscillator as a regionally stable ship-rolling
       model: dynamics, energy balance, and (regional) passivity}
\author{}
\date{}

\begin{document}
\maketitle

\noindent
This note collects the equations and the energy / passivity argument
behind the ship-rolling demo in
\texttt{python/notebooks/duffing/ship\_rolling\_simulation.py}.
Everything is stated rigorously enough to verify by hand. The dynamics
are taken verbatim from
\texttt{python/src/sysid/evaluation/true\_dynamics.py}, line~30:
\[
\ddot q \;=\; -\delta\,\dot q \;-\; q \;+\; q^3 \;+\; u,
\qquad \delta=0{.}3.
\]

\section{State-space form}

Let $x = (q, \dot q)^\top \in \R^2$. With input $u\in\R$ and damping
$\delta > 0$, the dimensionless softening Duffing oscillator reads
\begin{equation}
\dot x
  \;=\; f(x, u)
  \;=\;
  \begin{pmatrix} \dot q \\[2pt]
                  -\delta\,\dot q \;-\; q \;+\; q^3 \;+\; u
  \end{pmatrix}.
\label{eq:state}
\end{equation}
The cubic enters with a {\bfseries plus} sign --- this is the
{\itshape softening} Duffing, not the hardening one. The conservative
restoring force is therefore
\begin{equation}
F_{\mathrm{cons}}(q) \;=\; -q + q^3,
\label{eq:Fcons}
\end{equation}
which {\itshape weakens} as $|q|$ grows past $1$, vanishes at
$q=\pm 1$, and {\itshape reverses sign} for $|q|>1$.

\section{Mapping to ship rolling}

Let $\varphi$ denote the roll angle and $\varphi_v$ the
\emph{angle of vanishing stability} (the heel angle past which the GZ
righting arm changes sign --- physically, the capsize boundary). With
the dimensionless rescaling
\begin{equation}
q \;=\; \varphi / \varphi_v,
\qquad
\dot q \;=\; \dot\varphi / \varphi_v,
\label{eq:rescale}
\end{equation}
\eqref{eq:state} corresponds, after multiplying through by $\varphi_v$,
to the ship-roll model
\begin{equation}
\ddot\varphi
  \;=\; -\delta\,\dot\varphi
        \;-\; \omega_0^2\,\varphi
        \;+\; k_3\,\varphi^3
        \;+\; \tilde u,
\label{eq:phi-eqn}
\end{equation}
where $\omega_0^2$ and $k_3>0$ are obtained from the cubic
Taylor truncation of the GZ righting arm
\(
GZ(\varphi) \approx C_1\,\varphi + C_3\,\varphi^3,
\)
with $C_1>0$, $C_3<0$ (softening). The dimensionless
form~\eqref{eq:state} is~\eqref{eq:phi-eqn} with $\omega_0=1$, $k_3=1$,
and $\varphi_v$ absorbed into the state. The hilltops of the
dimensionless potential at $q=\pm 1$ correspond to
$\varphi = \pm \varphi_v$.

\section{Potential and Hamiltonian}

The conservative force~\eqref{eq:Fcons} is the negative gradient of
\begin{equation}
V(q) \;=\; \tfrac{1}{2}\,q^2 \;-\; \tfrac{1}{4}\,q^4,
\qquad
V'(q) \;=\; q - q^3 \;=\; -F_{\mathrm{cons}}(q).
\label{eq:V}
\end{equation}
Properties of $V$:
\begin{itemize}
\item Local minimum at $q=0$ with $V(0) = 0$ (upright equilibrium).
\item Local maxima (\emph{hilltops}) at $q=\pm 1$ with
      \(V(\pm 1) = \tfrac{1}{2} - \tfrac{1}{4} = \tfrac{1}{4}\).
\item Unbounded below: $V(q) \to -\infty$ as $|q|\to\infty$.
\end{itemize}
The total mechanical energy (which is the Hamiltonian when
$u = \delta = 0$) is
\begin{equation}
H(q,\dot q)
  \;=\; \underbrace{\tfrac{1}{2}\,\dot q^2}_{\text{kinetic}}
        \;+\; \underbrace{\tfrac{1}{2}\,q^2 - \tfrac{1}{4}\,q^4}_{V(q)}.
\label{eq:H}
\end{equation}

\paragraph{Code/text sign convention.} The function
\texttt{duffing\_V\_energy} in
\texttt{python/src/sysid/evaluation/true\_dynamics.py} returns
$V_{\mathrm{code}} := -H$ (sign-flipped) so that the rule
{\it``inside basin''} reads as \(V_{\mathrm{code}} > V_{\mathrm{saddle,code}} = -1/4\).
In the convention~\eqref{eq:H} the equivalent rule is $H < 1/4$.

\section{Equilibria and the saddle-node bifurcation in $u$}

For constant input $u$, equilibria of~\eqref{eq:state} satisfy
$\dot q = 0$ and
\begin{equation}
P_u(q) \;:=\; q^3 - q + u \;=\; 0.
\label{eq:cubic}
\end{equation}
The discriminant of~\eqref{eq:cubic} is $\Delta = 4 - 27\,u^2$. Hence:
\begin{itemize}
\item $|u| < u_c$: three real roots; one stable equilibrium near
      $q=0$ and two saddles.
\item $|u| = u_c$: a stable equilibrium and a saddle merge
      (saddle-node).
\item $|u| > u_c$: only one real root, an unstable equilibrium far
      from $q=0$. {\bfseries No stable equilibrium exists.}
\end{itemize}
The critical input is
\begin{equation}
u_c \;=\; \frac{2}{3\sqrt{3}} \;\approx\; 0.3849.
\label{eq:uc}
\end{equation}
For $u=0$, the three equilibria are $q\in\{-1,0,+1\}$. Linearising
$f(\cdot, 0)$ at $(0,0)$ gives Jacobian
\(
A_0 = \bigl(\begin{smallmatrix} 0 & 1 \\ -1 & -\delta \end{smallmatrix}\bigr)
\),
characteristic equation $\lambda^2 + \delta\lambda + 1 = 0$, eigenvalues
with negative real part for any $\delta > 0$ ⇒ asymptotically stable.
Linearising at $(\pm 1, 0)$ gives Jacobian
\(
A_{\pm} = \bigl(\begin{smallmatrix} 0 & 1 \\ 2 & -\delta \end{smallmatrix}\bigr)
\),
characteristic equation $\lambda^2 + \delta\lambda - 2 = 0$,
eigenvalues $\lambda_{1,2} = (-\delta \pm \sqrt{\delta^2 + 8})/2$:
one positive, one negative ⇒ saddle.

\section{Energy balance}

Multiply~\eqref{eq:state}'s second component by $\dot q$:
\begin{align}
\dot q\,\ddot q + \delta\,\dot q^2 + (q - q^3)\,\dot q
   &= u\,\dot q.
\notag\\
\frac{d}{dt}\!\left[\tfrac{1}{2}\dot q^2
       + \tfrac{1}{2}q^2 - \tfrac{1}{4}q^4\right]
   &= -\delta\,\dot q^2 \;+\; u\,\dot q.
\end{align}
That is,
\begin{equation}
\boxed{\;\;
\dot H \;=\; -\delta\,\dot q^2 \;+\; u\,\dot q
\;\;}
\label{eq:dHdt}
\end{equation}
along trajectories of~\eqref{eq:state}. Equation~\eqref{eq:dHdt}
is the central energy identity:
\begin{itemize}
\item $-\delta\,\dot q^2 \le 0$: damping continuously dissipates
      energy.
\item $u\,\dot q$: power injected by the input through the velocity
      port; sign depends on $\mathrm{sign}(u\,\dot q)$.
\end{itemize}

\section{Passivity --- but only \emph{regionally}}

Choose $u$ as the input and the collocated output
\(
y \;:=\; \dot q.
\)
Substituting in~\eqref{eq:dHdt},
\begin{equation}
\dot H
  \;=\; -\delta\,y^2 \;+\; u\,y
  \;\le\; u\,y.
\label{eq:passivity-ineq}
\end{equation}
This is the passivity differential inequality
$\dot S \le u^\top y$ with candidate storage $S = H$.
\textbf{However}, $H$ is {\bfseries not} bounded below on $\R^2$:
indeed $V(q) \to -\infty$ as $|q|\to\infty$. A valid storage function
for input--output passivity must be bounded below (and is conventionally
nonnegative with $S(0)=0$). We therefore have:

\paragraph{Regional passivity.}
On the open sublevel set
\begin{equation}
\Omega
  \;:=\;
  \bigl\{ (q,\dot q)\in\R^2 :\, H(q,\dot q) < \tfrac{1}{4} \bigr\}
  \;\cap\;
  \bigl\{|q| < 1\bigr\},
\label{eq:Omega}
\end{equation}
the function $H$ is nonnegative, $H(0,0)=0$, and
$\dot H \le u\,y$ holds. So~\eqref{eq:state} is
\emph{passive on $\Omega$} with storage $H$ and supply rate
$w(u,y)=u\,y$. The set $\Omega$ is the (open) basin of attraction of
the upright equilibrium under $u\equiv 0$, capped at the saddle level
$H = V(\pm 1) = 1/4$.

\paragraph{Failure of global passivity.}
As soon as the trajectory leaves $\Omega$ --- i.e.\ crosses the
hilltop ---
$H$ ceases to be sign-definite. The {\itshape inequality}
$\dot H \le u\,y$ still holds (it is an algebraic identity), but $H$
becomes arbitrarily negative because $V$ does, so $H$ no longer
certifies boundedness or stability. The Duffing system in this form is
{\bfseries not globally passive} with respect to the supply rate
$u\,\dot q$ and storage $H$.

\section{``Where does the unbounded kinetic energy come from?''}

Suppose $u\equiv 0$ and the trajectory has just barely crossed a
hilltop, with finite $H_0 := H(q_0, \dot q_0)$ at $t=0$,
$|q_0| > 1$. By~\eqref{eq:dHdt} the total mechanical energy is
non-increasing along the trajectory:
\begin{equation}
H(t) \;\le\; H_0 \quad\text{for all } t \ge 0.
\label{eq:Hbound}
\end{equation}
Substituting~\eqref{eq:H} and rearranging,
\begin{equation}
\tfrac{1}{2}\,\dot q(t)^2
  \;\le\;
  H_0 \;-\; \tfrac{1}{2}\,q(t)^2 \;+\; \tfrac{1}{4}\,q(t)^4,
\end{equation}
so that, asymptotically as $|q|\to\infty$,
\begin{equation}
\dot q^2 \;\sim\; \tfrac{1}{2}\,q^4,
\qquad
\text{i.e.}\qquad
\underbrace{\tfrac{1}{2}\dot q^2}_{\text{KE}}
  \;\sim\; \tfrac{q^4}{4}
  \;\to\; +\infty,
\quad
\underbrace{V(q)}_{\text{PE}}
  \;\sim\; -\tfrac{q^4}{4}
  \;\to\; -\infty.
\label{eq:asymp}
\end{equation}
The kinetic energy grows without bound, the potential energy drops
without bound, and they cancel to keep the {\itshape sum} $H$
bounded above by $H_0$ (in fact slowly decreasing because of damping).

\paragraph{The model does not create or destroy energy.}
What happens past the hilltop is exactly what happens to a marble
rolling off the side of a hill into a bottomless pit: kinetic energy
gained equals potential energy lost, instant by instant, with damping
removing a small amount as heat. The {\itshape source} of the
``infinite kinetic energy'' is the polynomial potential, which is
unbounded below {\bfseries by construction} --- it is a low-order
Taylor expansion of a real, finite restoring law (e.g.\ a GZ curve)
that has been extrapolated outside its domain of validity.

\paragraph{Physical interpretation in the ship-rolling picture.}
For a real ship the GZ curve flattens out and eventually changes sign
again past the angle of vanishing stability (the hull has a second,
fully inverted, equilibrium --- the ``turtle'' position). The true
potential is bounded below. The Duffing model is therefore a faithful
description of the rolling dynamics {\itshape within the basin of
attraction of upright}, and its divergence is the sharp signature of
\emph{leaving} that basin --- equivalent to capsize. The model is
silent about what posture the ship ends up in (turtle, sunk, broken),
which is fine: the regional verification question is precisely
``does the trajectory stay in the basin?'', not ``what happens
afterwards?''.

\section{Implications for regional verification}

The regional analogue of global $L_2$-stability or global passivity
is, informally:

\begin{quote}
For all admissible input signals $u(\cdot)\in\mathcal U$ and all
initial conditions $x_0 \in X_0 \subseteq \Omega$, the resulting
trajectory $x(\cdot)$ remains in some forward-invariant set
$\Omega' \subseteq \Omega$, equivalently, the storage $H(x(t))$
remains bounded above by some $\overline H < \tfrac{1}{4}$ for all
$t \ge 0$.
\end{quote}

This is the natural specification class for benchmarks like the
softening Duffing where global guarantees are impossible by
construction, and it is the property that an identified model with
``regionally stable'' structural constraints is meant to certify.

\appendix
\section{Key numerical values used in the demo}

\begin{tabular}{ll}
\hline
Symbol & Value (default in the demo) \\
\hline
$\delta$ (damping)              & $0.3$ \\
Sample time $T_s$               & $0.05$ s \\
Saddle-node threshold $u_c$     & $2/(3\sqrt{3}) \approx 0.3849$ \\
Saddle level of $H$             & $1/4$ \\
Angle of vanishing stability $\varphi_v$ (visual scaling)
                                & $60^\circ$ (CLI flag \verb|--phi-v-deg|) \\
Stable forcing $u(t)$           & $0.25\,\sin(0.4\,t)$ \\
Capsize forcing $u(t)$          & $1.5$ on $t\in[4,6]$ s, else $0$ \\
\hline
\end{tabular}

\end{document}
